\documentclass[aspectratio=169,10pt,spanish]{beamer}

\input{../../../_preambulo_beamer.tex}
\input{../../../_comandos_beamer.tex}

\selectlanguage{spanish}

\title{MA1001B - Modelación Estadística para la Toma de Decisiones}
\subtitle{Lección 13: Variables aleatorias discretas: FMP, FDA, esperanza, varianza}
\author{}
\institute{Tecnol\'ogico de Monterrey}
\date{}

\begin{document}

\begin{frame}[plain]
  \vspace{0.9cm}
  {\Large\bfseries MA1001B - Modelación Estadística para la Toma de Decisiones\par}
  \vspace{0.55cm}
  {\large Lección 13: Variables aleatorias discretas\par}
  \vspace{0.15cm}
  {\large FMP, FDA, esperanza, varianza\par}
  \vspace{0.55cm}
  {\color{TecRojo}\rule{\textwidth}{1.2pt}}
  \vspace{0.6cm}

  Facultad MA1001B\\[0.3cm]
  Periodo\\[0.3cm]
  Tecnol\'ogico de Monterrey
\end{frame}

\begin{frame}{Un conteo de defectos como variable aleatoria}
  \begin{alertblock}{Pregunta guía}
    ¿Cómo se especifica una variable aleatoria discreta, y por qué $E[X^2]$
    no es el mismo número que $(E[X])^2$?
  \end{alertblock}

  \vspace{0.6em}
  Al final de esta lección, usted puede:
  \begin{itemize}
    \item comprobar que una FMP tiene masas no negativas que suman 1;
    \item acumular la FMP en una FDA;
    \item calcular $E[X]$ y $E[X^2]$ como sumas ponderadas;
    \item obtener $\Var(X)=E[X^2]-(E[X])^2$.
  \end{itemize}

  \vfill
  \scriptsize Todos los conteos de defectos usados hoy son sintéticos. No se usan datos reales de empresa.
\end{frame}

\begin{frame}{Soporte y masas de probabilidad}
  \small
  Sea $X$ el conteo de defectos en un lote sintético, $X\in\{0,1,2,3\}$.
  \[
    P(X=0)=0.50,\quad
    P(X=1)=0.30,\quad
    P(X=2)=0.15,\quad
    P(X=3)=0.05.
  \]
  Una función $p$ es una FMP cuando $p(x)\ge 0$ para todo $x$ y
  $\sum_x p(x)=1$.
\end{frame}

\begin{frame}[fragile]{Paso 1: Función de masa de probabilidad}
  \begin{columns}[T]
    \begin{column}{0.42\textwidth}
      \small
      \begin{lstlisting}
support = [0, 1, 2, 3]
pmf = [0.50, 0.30, 0.15, 0.05]
mass_sum = sum(pmf)
      \end{lstlisting}

      \begin{block}{Salida verificada}
        $P(X=0)=0.500000$\\
        $P(X=1)=0.300000$\\
        $P(X=2)=0.150000$\\
        Suma $=1.000000$
      \end{block}
    \end{column}
    \begin{column}{0.58\textwidth}
      \centering
      \includegraphics[width=\textwidth]{../figuras/va_discreta_01_fmp.png}
      \vspace{0.2em}
      \scriptsize Gráfica de tallos de la FMP del Paso 1
    \end{column}
  \end{columns}
\end{frame}

\begin{frame}{La FDA acumula la FMP}
  \small
  \[
    F(x)=P(X\le x)=\sum_{t\le x}P(X=t).
  \]
  Valores verificados:
  \[
    F(0)=0.50,\quad F(1)=0.80,\quad F(2)=0.95,\quad F(3)=1.00.
  \]
  Las probabilidades de intervalo vienen de diferencias:
  \[
    P(1<X\le 3)=F(3)-F(1)=0.20.
  \]
\end{frame}

\begin{frame}[fragile]{Paso 2: Función de distribución acumulada}
  \begin{columns}[T]
    \begin{column}{0.40\textwidth}
      \small
      \begin{lstlisting}
cdf = np.cumsum(pmf)
p_hi = cdf[3] - cdf[1]
      \end{lstlisting}

      \begin{block}{Salida verificada}
        $F(1)=0.800000$\\
        $F(3)=1.000000$\\
        $P(1<X\le 3)=0.200000$
      \end{block}
    \end{column}
    \begin{column}{0.60\textwidth}
      \centering
      \includegraphics[width=\textwidth]{../figuras/va_discreta_02_fda.png}
      \vspace{0.2em}
      \scriptsize FDA de función escalón del Paso 2
    \end{column}
  \end{columns}
\end{frame}

\begin{frame}{Paso 3: $E[X^2]$ no es $(E[X])^2$}
  \begin{columns}[T]
    \begin{column}{0.42\textwidth}
      \small
      \[
        E[X]=0.750000,
      \]
      \[
        (E[X])^2=0.562500,
      \]
      \[
        E[X^2]=1.350000,
      \]
      \[
        \Var(X)=0.787500.
      \]

      \vspace{0.3em}
      \textbf{Límite:} $E[X^2]\ne(E[X])^2$. La Lección 14 especializa este
      lenguaje a ensayos de Bernoulli y al modelo binomial.
    \end{column}
    \begin{column}{0.58\textwidth}
      \centering
      \includegraphics[width=\textwidth]{../figuras/va_discreta_03_esperanza_varianza.png}
      \vspace{0.2em}
      \scriptsize Momentos y varianza del Paso 3
    \end{column}
  \end{columns}
\end{frame}

\begin{frame}{Verificación de conceptos}
  \begin{enumerate}
    \item ¿Por qué $P(X=4)=0$ para esta $X$?
    \item Calcule $F(2)$ a partir de la FMP.
    \item Calcule $E[X]$ a partir de las cuatro masas.
    \item ¿Por qué $1.350000$ no es igual a $0.562500$?
  \end{enumerate}

  \vspace{0.6em}
  \begin{block}{Conexión con la Lección 14}
    La sesión puede continuar directamente con ensayos de Bernoulli, la FMP
    binomial y un supuesto de probabilidad de éxito constante.
  \end{block}

  \vfill
  \scriptsize Fuente: OpenStax, \textit{Introductory Statistics 2e},
  Secciones 4.1--4.2. CC BY-NC-SA.
\end{frame}

\appendix



\end{document}
